\documentclass{article}
\begin{document}
\section*{Wolf Proposition 6.1 --- Russo--Dye constant}

\paragraph{Source statement.}
Wolf, \textit{Quantum Channels \& Operations}, Proposition~6.1:
``If $T$ is a positive map on $\mathcal M_d(\mathbb C)$, then
$\varrho(T)\le \|T(\mathbf 1)\|_\infty$.''

\paragraph{Deviation.}
The source proof invokes the Russo--Dye theorem $\|T(X)\|_\infty\le \|T(\mathbf 1)\|_\infty\,\|X\|_\infty$
to bound every eigenvalue by $\|T(\mathbf 1)\|_\infty$ and thus the spectral
radius.  The present formalization (Mathlib) provides
\texttt{PositiveLinearMap.norm\_apply\_le\_of\_nonneg} for PSD inputs
and the four-part decomposition for general inputs, which yields the weaker
bound $\varrho(T)\le 4\,\|T(\mathbf 1)\|_\infty$ (factor~$4$ instead of~$1$).

The trace-preserving and unital cases circumvent this because the eigenvalue
bound $|\lambda|\le 1$ is obtained directly from the orbit argument in
\texttt{IsPositiveMap.eigenvalue\_norm\_le\_one\_of\_tracePreserving}
(\texttt{TNLean/Channel/Determinant/Bound.lean}) and eigenvalue~$1$ existence
from the Perron--Frobenius theorem (\texttt{TNLean/Channel/Peripheral/SpectralRadius.lean}).

\paragraph{Elimination plan.}
Formalize the Russo--Dye theorem for matrix algebras (either by importing the
general C$^\ast$-algebra result from Mathlib once available, or by a direct
matrix argument using the Russo--Dye theorem for the special case of
$\mathcal B(\mathcal H)$).  Until then the trace-preserving/unital cases
are fully covered; the general positive-map case remains with factor~$4$.

\paragraph{Tracking.} Issue \#3984.
\end{document}
